\homework{1}{Sun 21 Aug 2022 23:38}{Week 1 due Aug 30}

\begin{enumerate}
	\item 4.H. If $\xi \in \mathbf{R}$ is irrational and $r \in \mathbf{R}, r \neq 0$, is rational, show that $r+\xi$ and $r \xi$ are irrational.
\begin{proof}
	We use, without proving, the fact that the rational numbers $\mathbb{Q}$ forms a field. In particular, it is closed under multiplication (and division, by existence of multiplicative inverse) and addition (and substraction, by existence of additive inverse).

We prove the contrapositive. Suppose $r+\xi$ is not irrational. Then $r + \xi $ is rational. Because $r $ is rational, $\xi = \left( r + \xi  \right) -r$ is rational, which means it's not irrational.

The same proof for $r \xi $. Modify the previous proof by exchanging addition and multiplication (substraction and division).
\end{proof}
	\item 5.A. If $a, b \in \mathbf{R}$ and $a^{2}+b^{2}=0$, show that $a=b=0$.
\begin{proof}
Suppose for the sake of contradiction that either $a\neq 0$ or $b\neq 0$. By symmetry, it suffices to only prove the case for $a \neq 0$. By Theorem 5.5(a), $a^2 >0$. 

Now we want to prove $b^2 \ge 0$. If $b = 0$, it follows from Theorem 4.5(a) that $b^2 = 0 \cdot 0 = 0$. If $b \neq  0$, then $b^2 > 0$ by Theorem 5.5(a).

In the first case, $a^2+b^2 = a^2 + 0 = a^2 > 0$. In the second case, $a^2 + b^2 > 0 + 0 = 0$ by Theorem 5.6(b). In both cases $a^2 + b^2 > 0$, but it is assumed that $a^2 + b^2 = 0$.
\end{proof}
	\item 5.C. If $a>-1, a \in \mathbf{R}$, show that $(1+a)^{n} \geq 1+n a$ for all $n \in \mathbf{N}$. This inequality is called Bernoulli's Inequality. (Hint: use mathematical induction.)
\begin{proof}
	Induct on $n$. 

	Base case: We need to show $(1+a)^0 \ge 1+0 \cdot a$. This simplifies to  $1 \ge 1$ which is trivial.

	Induction step: We need to show $(1+a)^{n+1}\ge 1+(n+1)a $.
	\begin{align*}
		(1+a)^{n+1} &= (1+a)^{n}(1+a)\\
		\intertext{By induction hypothesis, $(1+a)^n \ge 1+na$. Apply Theorem 5.6(c) (note that $1+a>0$):}
			 &\ge (1+na)(1+a)\\
			 &= 1+a+na+na^2
		\intertext{Since $a^2\ge 0$:}
			 &\ge  1+(n+1)a 
	.\end{align*}
	
\end{proof}

	\item 5.O. If $x, y, z$ belong to ${\mathbf{R}}$, then $x \leq y \leq z$ {\color{red} or $z\le y\le x$} if and only if $|x-y|+|y-z|=|x-z|$.
\begin{proof}
	We split the goal into two subgoals, which together imply the goal:
	\[
	\begin{cases}
		z\le y\le x \iff |x-y|+|y-z|=|x-z| \land x\ge z\\
		x\le y\le z \iff |x-y|+|y-z|=|x-z| \land x\le  z
	\end{cases}
	.\] 
	By symmetry, we only need to prove the first subgoal. 

	 $\implies$:
	Since $y\ge z$, $y-z \ge 0$, hence $|y-z| = y-z$. Similarly $|x-y| = x-y $ and $|x-z|=x-z $. The goal now simplifies to $(x-y)+(y-z)=x-z$, which follows by simplification.

	$\impliedby$:
	It suffices to show that both $y>x $ and $y<z $ cannot be true. By symmetry, we only show that $y>x $ is false.
	
	Suppose for the sake of contradiction that $y>x$. By similar reasoning, $|x-y|+|y-z| = (y-z)+(y-z)$. By assumption, $x\ge z$, so $|x-z| = x-z$. Since $y>x$, $y-z>x-z$. Therefore, $|x-y|+|y-z| > (y-z)+(x-z) = (y-z)+|x-z| > |x-z|$, a contradiction. 

\end{proof}
	\item 6.D. Give an example of a set of irrational numbers that has a rational supremum.
\begin{proof}
	Consider the set \[
	\{-\sqrt{2}/n:n \in \mathbb{N} \} 
	.\] 
	
\end{proof}
	\item 6.G. If $S$ is a bounded set in $\boldsymbol{R}$ and if $S_{0}$ is a nonempty subset of $S$, then show that
\[
\inf S \leq \inf S_{0} \leq \sup S_{0} \leq \sup S .
\]
Sometimes it is more convenient to express this in another way. Let $D \neq \emptyset$ and let $f: D \rightarrow \boldsymbol{R}$ have bounded range. If $D_{0}$ is a non-empty subset of $D$, then
\[
\inf \{f(x): x \in D\} \leq \inf \left\{f(x): x \in D_{0}\right\} \leq \sup \left\{f(x): x \in D_{0}\right\} \leq \sup \{f(x): x \in D\} .
\]
\begin{proof}
	The second inequality follows from definition, so we only need to prove the first and third inequalities. By symmetry, we only need to prove the first inequality.

	It suffices to show that $\operatorname{inf}S$ is a lower bound of $S_0$. By definition, for all $x \in  S$, $x\ge \operatorname{inf}S$. Take $s \in  S_0$. Since $S_0 \subset S$, $s \in  S$. Therefore, $s\ge \operatorname{inf}S$.
\end{proof}
	\item 6.H. Let $X$ and $Y$ be non-empty sets and let $f: X \times Y \rightarrow R$ have bounded range in $\boldsymbol{R}$. Let
\[
f_{1}(x)=\sup \{f(x, y): y \in Y\}, \quad f_{2}(y)=\sup \{f(x, y): x \in X\} .
\]
Establish the Principle of Iterated Suprema:
\[
\begin{aligned}
\sup \{f(x, y): x \in X, y \in Y\} &=\sup \left\{f_{1}(x): x \in X\right\} \\
&=\sup \left\{f_{2}(y): y \in Y\right\} .
\end{aligned}
\]
We sometimes express this in symbols by:
\[
\sup _{x, y} f(x, y)=\sup _{x} \sup _{y} f(x, y)=\sup _{y} \sup _{x} f(x, y) .
\]
\begin{proof}
	By symmetry, it suffices to show \[
	\sup \left\{ f(x, y): x \in X, y \in Y\right\}=\sup \left\{f_{1}(x): x \in X\right\}
. \]
	We first show $\sup \left\{ f(x, y): x \in X, y \in Y\right\}\le \sup \left\{f_{1}(x): x \in X\right\}$.
	To show this, we need to show that $ \sup \left\{f_{1}(x): x \in X\right\}$ is an upper bound of $ \left\{ f(x, y): x \in X, y \in Y\right\}$. For any $(x_0,y_0) \in  X \times Y$, $f_1(x_0) = \sup \left\{ f(x_0,y):y \in  Y \right\} $. Now $f(x_0,y_0) \in \left\{ f(x_0,y):y \in  Y \right\}$, therefore $\sup \left\{ f(x_0,y):y \in  Y \right\} \ge f(x_0,y_0)$.

	Next we show $\sup \left\{ f(x, y): x \in X, y \in Y\right\}\ge \sup \left\{f_{1}(x): x \in X\right\}$.
	To show this, we need to show that $\sup \left\{ f(x, y): x \in X, y \in Y\right\}$ is an upper bound of $\left\{f_{1}(x): x \in X\right\}$.
	Take  $x_0 \in  X$. Now $f_1(x_0) = \sup \left\{ f(x,y): x,y \in \{x_0 \} \times Y \right\} \le \sup \left\{ f(x,y): x,y \in X \times Y \right\} $, where the inequality comes from exercise 6G.
\end{proof}
\end{enumerate}
